\chapter{The Shifted Boundary Method}
\label{ch:c4}

\emph{Module: \file{04\_shifted\_boundary/}. Flow past a body whose
boundary is OFF the grid; the Taylor shift $d\neq0$ does real work;
surrogate-traction drag.}

\section{Surrogate vs true boundary}
The true immersed boundary $\Gamma$ (a level set of an oracle --- a
\code{Box} or \code{Sphere}) generally cuts through octree cells. SBM does
\emph{not} body-fit. Instead it classifies cells and keeps a
\textbf{surrogate boundary} $\Gambar$ made of \emph{whole} element faces
(\file{sbm/surrogate.py:14--21}):
\begin{itemize}
\item \code{classify_lambda(tree, oracle, lam, domain)} marks a cut cell as
  \emph{retained} iff its domain-\emph{outside} volume fraction $\le\lambda$.
  So $\lambda=0$ keeps only fully-interior cells (surrogate strictly inside
  $\Omega$; if the body is cell-aligned the surrogate coincides with
  $\Gamma$, $d=0$); $\lambda=0.5$ is the flow-past value.
\item \code{extract_surrogate(tree)} returns the surrogate faces: faces of
  retained cells with no retained neighbour across them and not on the outer
  domain boundary (whole-face invariant).
\end{itemize}

\section{The Taylor shift $S = N + \grad N\cdot d$}
Let $d$ be the vector from a surrogate Gauss point to its closest point on
$\Gamma$. The shifted trial/test operator is a first-order Taylor
extrapolation from the surrogate to the true boundary:
\begin{equation}
S\,N_a = N_a + (\grad N_a)\cdot d.
\end{equation}
The Dirichlet data is imposed at the mapped point $y=\text{xq}+d$, and the
shifted trial value is $Su = u + (\grad u)\cdot d$. When $d=0$
(cell-aligned body), $S N_a = N_a$ and SBM is exactly standard Nitsche ---
the $d=0$ case of Ch.~\ref{ch:c3}.

\begin{keybox}\textbf{Key idea.} SBM is Nitsche with the shape functions
Taylor-extrapolated to the true boundary. No re-meshing, no cut-cell
quadrature: the shift is fixture \emph{data} ($d$, $\text{corr}$), not
solver code. The same predictor/PPE machinery runs
unchanged.\end{keybox}

\section{Surrogate normal, distance, area correction}
\code{GeometryData.evaluate} (\file{sbm/surrogate.py:222--253}) produces the
frozen FP64 per-epoch geometry cache at surrogate Gauss points:
\begin{itemize}
\item \textbf{$d$} --- the shift vector (surrogate GP $\to$ closest point on
  $\Gamma$), from the oracle's \code{distance_vector}.
\item \textbf{$n$} --- the \emph{true-boundary} outward normal
  $n=\pm\grad\psi/|\grad\psi|$, always pointing \emph{out of} the
  computational domain $\Omega$.
\item \textbf{$\text{corr}$} --- the area correction
  $\text{corr}=\tilde n\cdot n$, where $\tilde n$ is the surrogate face's
  grid-aligned unit normal. It converts a surrogate-face measure to the true
  boundary measure ($\text{corr}=1$ when $\Gambar=\Gamma$).
\end{itemize}

\section{Force / drag on the immersed body}
The drag reads the traction on the surrogate faces, area-corrected to the
true boundary. \code{surrogate_traction(dm, sf, geo, x_all, nu, ndof)}
(\file{sbm/vector.py:562--605}) computes
\begin{equation}
F = \oint_{\Gambar}\big(p\,n - \nu\,(\grad u)^\top n\big)\,\text{corr}\,\mathrm{d}\Gambar,
\end{equation}
with the crucial \textbf{orientation contract}: \code{geo.n} is
domain-outward (= \emph{into} the obstacle for exterior flow), so the
obstacle-outward normal is $\hat n=-\texttt{geo.n}$ and the integrand is
$+p\,\texttt{geo.n}-\nu(\grad u)^\top\texttt{geo.n}$. At $d=0$ this reduces to
the raw body-fitted traction.

\section{Verification: a genuine sub-cell shift}
The same square as Ch.~\ref{ch:c3}, but its center is offset by a sub-cell
fraction ($\text{offset}=0.05$), so the grid-aligned surrogate no longer
coincides with the true \code{Box} face: $0<d_{\max}<h$. Now the Taylor
$(\grad N_a)\cdot d$ term \emph{and} the area correction $\text{corr}$ do
real work. The stepper setup is \emph{identical} to the weak path of
Ch.~\ref{ch:c3} --- the shift is fixture data. The monolithic oracle carries
the \emph{same} shifted geometry, so it is the same-mesh-with-shift oracle.

At level~4, $\Rey=40$, $\text{offset}=0.05$:
\[
d_{\max}=\ShDmax\quad(0<d_{\max}<h=0.0625,\ d_{\max}/h=\ShDmaxRatio)
\]
\begin{center}
\begin{tabular}{@{}lccc@{}}
\toprule
 & projection & monolithic (true shift) & rel-diff\\
\midrule
$\Cd$ & $+$\ShCdProj & $+$\ShCdMono & \textbf{\ShCdRel}\\
\bottomrule
\end{tabular}
\end{center}
The consistent-projection split \emph{with the genuine SBM shift active}
matches the same-mesh-with-shift monolithic to \ShCdRel{} in drag ---
\textbf{projection $+$ SBM works end-to-end in 2-D}.

\begin{keybox}\textbf{Anti-vacuity (the shift is load-bearing).} Run with
\code{--zero-shift} to set $\texttt{geo.d}=0$, $\texttt{geo.corr}=1$ while
keeping everything else fixed; the match against the TRUE shifted oracle
breaks --- proof the Taylor term and area correction are doing real
work.\end{keybox}

\begin{runbox}\textbf{Run it.} \code{python run.py} in
\file{04\_shifted\_boundary/} marches the shifted projection vs the shifted
monolithic. \code{--zero-shift} demonstrates the anti-vacuity break;
\code{--offset 0} returns to the body-fitted $d=0$ case.\end{runbox}

\question At $\text{offset}=0.05$, $d_{\max}/h=0.80$: the shift is nearly a
full cell. What happens to the Taylor approximation as $\text{offset}\to h$?
Why must the offset stay $<h$?

\question Only a handful of surrogate Gauss points are area-corrected
($\text{corr}\neq1$). Which ones, geometrically, and why not all of them?
